\documentclass[a4paper,nopdfoutputerror]{quantumarticle}
\pdfoutput=1

\usepackage[utf8]{inputenc}
\usepackage[T1]{fontenc}
\usepackage{amsmath,amssymb,amsthm}
\usepackage{graphicx}
\usepackage{booktabs}
\usepackage{xcolor}
\usepackage{algorithm}
\usepackage{algpseudocode}
\usepackage{microtype}
\usepackage{subcaption}
\usepackage{colortbl}
\usepackage{hyperref}
\usepackage[numbers,sort&compress]{natbib}

\definecolor{qnggreen}{RGB}{46,125,50}
\definecolor{euclidred}{RGB}{239,83,80}
\definecolor{classblue}{RGB}{13,71,161}

\title{Quantum Natural Gradient as the Inner Loop of Model-Agnostic Meta-Learning
for Variational Quantum Circuits}

\author{Daniel Mart\'inez P\'erez}
\affiliation{Quantum Machine Learning Group, Universidad Pablo de Olavide, Sevilla, Spain}
\email{dmarper2@upo.es}

\begin{document}
\maketitle

\begin{abstract}
Few-shot learning---the ability to generalize from only a handful of
examples---remains a central challenge in quantum machine learning.
Model-Agnostic Meta-Learning (MAML) addresses this by learning an initial
parameter configuration that can be rapidly adapted to new tasks via gradient
descent; yet when applied to Variational Quantum Circuits (VQCs), the
Euclidean gradient of the inner loop ignores the non-Euclidean geometry of
the quantum parameter space, leading to slow adaptation and brittleness under
hardware noise.
We propose \emph{Quantum Model-Agnostic Meta-Learning} (QMAML), a hybrid
quantum-classical framework that replaces the Euclidean inner loop of MAML
with a \emph{Quantum Natural Gradient} (QNG) step conditioned on the diagonal
of the Quantum Fisher Information Matrix (QFIM).
Because the QFIM is accessible exactly and efficiently via the parameter-shift
rule---at cost $\mathcal{O}(p)$ circuit evaluations, compared to
$\mathcal{O}(p^2)$ for the full classical Fisher matrix---this
geometric correction comes at manageable overhead while providing inner-loop
updates that respect the intrinsic curvature of the quantum state manifold.
Our architecture couples a classical encoder, shared and task-specific
\texttt{StronglyEntanglingLayers} in a 6-qubit VQC, and a linear classifier,
trained end-to-end via second-order MAML on the Omniglot benchmark in the
5-way $\{1,5\}$-shot setting.
Experimental results demonstrate that QMAML-QNG converges faster during
meta-training and achieves superior few-shot accuracy compared to
QMAML with Euclidean inner loop and classical MAML, with the gap widening
further when zero-error noise mitigation is applied on real IBM Quantum hardware.
\end{abstract}


% ─── Introduction ──────────────────────────────────────────────────────────────
\section{Introduction}
\label{sec:intro}

The ability to learn new concepts from a small number of examples is a
hallmark of human cognition.
In classical machine learning, \emph{meta-learning}---or ``learning to
learn''---has emerged as a principled framework for endowing models with this
capability~\cite{finn2017maml,vinyals2016matching,snell2017prototypical}.
Among meta-learning algorithms, Model-Agnostic Meta-Learning
(MAML)~\cite{finn2017maml} stands out for its simplicity and generality: it
seeks an initialization of model parameters such that a few gradient steps on
a new task yield strong performance.

Variational quantum circuits (VQCs) have attracted considerable attention as
quantum analogues of neural networks~\cite{cerezo2021variational,biamonte2017}.
In the NISQ era~\cite{preskill2018quantum}, VQCs represent the most
accessible path to potential quantum advantage in learning tasks.
However, training VQCs poses distinctive challenges: barren
plateaus~\cite{mcclean2018barren} cause gradient variances to decay
exponentially with circuit width, and standard Euclidean gradient descent
fails to exploit the inherent geometric structure of the quantum parameter
space.

The \emph{Quantum Fisher Information Matrix}
(QFIM)~\cite{stokes2020quantum,meyer2021fisher} provides the natural
Riemannian metric on the quantum state manifold.
Quantum Natural Gradient (QNG)~\cite{stokes2020quantum} preconditions
gradients with the inverse QFIM, yielding updates that are invariant to
re-parameterization and empirically converge faster on single-task VQC
training.
A key advantage over the classical setting is that the diagonal of the QFIM
is accessible exactly via the parameter-shift rule at $\mathcal{O}(p)$ circuit
evaluations, where $p$ is the number of parameters.

In this work we ask: \emph{does replacing the Euclidean inner loop of MAML
with QNG accelerate few-shot adaptation in VQCs?}
We answer affirmatively through the QMAML framework, which integrates
QFIM-preconditioned inner-loop updates into the MAML outer loop.
Our main contributions are:

\begin{enumerate}
  \item We formalize QMAML, the first meta-learning algorithm whose inner
        loop uses the Quantum Natural Gradient for VQCs.
  \item We derive an efficient diagonal QFIM estimator via parameter-shift
        and show it adds only $\mathcal{O}(p)$ overhead per inner step.
  \item We demonstrate empirically on Omniglot 5-way $\{1,5\}$-shot that
        QMAML-QNG converges faster and generalizes better than both
        QMAML-Euclidean and classical MAML baselines.
  \item We provide an IBM Quantum Runtime integration enabling direct
        execution on NISQ hardware with zero-error-noise-extrapolation (ZNE)
        mitigation.
\end{enumerate}

% ─── Background ────────────────────────────────────────────────────────────────
\section{Background}
\label{sec:background}

\subsection{Model-Agnostic Meta-Learning}

MAML~\cite{finn2017maml} considers a distribution over tasks
$p(\mathcal{T})$.
Each task $\mathcal{T}_i$ consists of a support set
$\mathcal{D}_i^{tr} = \{(x_j, y_j)\}_{j=1}^{N \cdot K}$
(N classes, K shots) and a query set $\mathcal{D}_i^{test}$.
Given model parameters $\boldsymbol{\theta}$, the \emph{inner loop} adapts
to task $\mathcal{T}_i$ via a small number of gradient steps:
\begin{equation}
  \boldsymbol{\theta}'_i = \boldsymbol{\theta}
    - \alpha \nabla_{\boldsymbol{\theta}} \mathcal{L}_{\mathcal{T}_i}
      (\boldsymbol{\theta}; \mathcal{D}_i^{tr}),
\end{equation}
and the \emph{outer loop} (meta-update) minimizes the expected query loss
over adapted parameters:
\begin{equation}
  \min_{\boldsymbol{\theta}} \sum_{\mathcal{T}_i \sim p(\mathcal{T})}
    \mathcal{L}_{\mathcal{T}_i}(\boldsymbol{\theta}'_i;\, \mathcal{D}_i^{test}).
\end{equation}

\subsection{Variational Quantum Circuits}

A VQC is a parameterized quantum circuit $U(\boldsymbol{\theta})$ applied
to a fixed initial state $|0\rangle^{\otimes n}$.
Given a feature vector $\mathbf{x}$, an encoding unitary
$S(\mathbf{x})$ embeds the classical data into the Hilbert space, and
the circuit produces expectation values
\begin{equation}
  f_i(\mathbf{x}, \boldsymbol{\theta}) = \langle 0 |
    U^\dagger(\boldsymbol{\theta})\, S^\dagger(\mathbf{x})\,
    O_i\,
    S(\mathbf{x})\, U(\boldsymbol{\theta})
  | 0 \rangle,
\end{equation}
where $O_i$ are observables (here, $O_i = Z_i$).
Gradients are computed via the parameter-shift
rule~\cite{mitarai2018quantum}:
\begin{equation}
  \frac{\partial f}{\partial \theta_k}
    = \frac{f(\theta_k + \tfrac{\pi}{2}) - f(\theta_k - \tfrac{\pi}{2})}{2}.
\end{equation}

\subsection{Quantum Fisher Information Matrix}

The QFIM $\mathbf{F}(\boldsymbol{\theta})$ is the Fubini-Study metric
on the space of quantum states~\cite{stokes2020quantum,meyer2021fisher}:
\begin{equation}
  F_{ij}(\boldsymbol{\theta})
    = 4\,\mathrm{Re}\!\left[
        \langle \partial_i \psi | \partial_j \psi \rangle
        - \langle \partial_i \psi | \psi \rangle
          \langle \psi | \partial_j \psi \rangle
      \right],
\end{equation}
where $|\psi(\boldsymbol{\theta})\rangle = U(\boldsymbol{\theta})|0\rangle$.
The Quantum Natural Gradient preconditions the gradient:
\begin{equation}
  \boldsymbol{\theta}' = \boldsymbol{\theta}
    - \alpha\, \mathbf{F}^{+}(\boldsymbol{\theta})\,
      \nabla_{\boldsymbol{\theta}} \mathcal{L},
\end{equation}
where $\mathbf{F}^{+}$ denotes the Moore-Penrose pseudoinverse.
We approximate $\mathbf{F}$ with its diagonal, which can be estimated as:
\begin{equation}
  \hat{F}_{ii}(\boldsymbol{\theta})
    = \mathbb{E}_{\mathbf{x}}\!\left[
        \left(\frac{\partial \langle O \rangle}{\partial \theta_i}\right)^{\!2}
      \right],
  \label{eq:qfim_diag}
\end{equation}
computable via $2p$ parameter-shift evaluations.

% ─── Method ────────────────────────────────────────────────────────────────────
\section{QMAML: Quantum MAML with Natural Gradient}
\label{sec:method}

\subsection{Architecture}

Our model $f_{\boldsymbol{\Theta}}$ is a hybrid quantum-classical pipeline:
\begin{equation}
  \mathbf{x} \xrightarrow{\phi_{\mathrm{enc}}}
  \mathbf{z} \in \mathbb{R}^{n_q} \xrightarrow{S(\mathbf{z})}
  \text{VQC}(\boldsymbol{\theta}_{\mathrm{sh}}, \boldsymbol{\theta}_{\mathrm{tk}})
  \xrightarrow{W_{\mathrm{cls}}} \hat{y},
\end{equation}
where:
\begin{itemize}
  \item $\phi_{\mathrm{enc}}: \mathbb{R}^{d} \to [-1,1]^{n_q}$ is a
        two-layer MLP with $\tanh$ output ($d=784$, $n_q = 6$);
  \item $S(\mathbf{z}) = \bigotimes_{i=1}^{n_q} R_Y(\pi z_i)$ is
        angle encoding;
  \item $U(\boldsymbol{\theta}_{\mathrm{sh}}, \boldsymbol{\theta}_{\mathrm{tk}})$
        consists of $L_s = 2$ \emph{shared} \texttt{StronglyEntanglingLayers}
        followed by $L_t = 1$ \emph{task-specific}
        \texttt{StronglyEntanglingLayers};
  \item observables $\{Z_i\}_{i=0}^{n_q-1}$ yield a real-valued
        feature vector fed to a linear classifier $W_{\mathrm{cls}}$.
\end{itemize}

The parameter split is central to the meta-learning interpretation:
$\boldsymbol{\theta}_{\mathrm{sh}}$ (shared) encode task-agnostic
quantum representations updated in the outer loop, while
$\boldsymbol{\theta}_{\mathrm{tk}}$ (task) are adapted in the inner loop
for each new task.

\subsection{QFIM-Preconditioned Inner Loop}

For each episode, given support set
$\mathcal{D}^{tr} = \{(\mathbf{x}_j, y_j)\}$, the QNG inner loop is:

\begin{algorithm}[H]
\caption{QMAML Inner Loop (QNG)}
\label{alg:inner}
\begin{algorithmic}[1]
  \Require model $f$, support set $\mathcal{D}^{tr}$, steps $T$,
           learning rate $\alpha$, regularization $\varepsilon$
  \State $\mathbf{z} \gets \phi_{\mathrm{enc}}(\mathbf{x})$
         \quad \Comment{encode support features}
  \State $\hat{\mathbf{F}} \gets \mathrm{diag\text{-}QFIM}(\mathbf{z},\,
         \boldsymbol{\theta}_{\mathrm{sh}},\, \boldsymbol{\theta}_{\mathrm{tk}})$
         \quad \Comment{Eq.~\eqref{eq:qfim_diag}, parameter-shift}
  \State $\mathbf{G} \gets |\hat{\mathbf{F}}| + \varepsilon$
         \quad \Comment{regularized diagonal metric}
  \For{$t = 1, \ldots, T$}
    \State $\hat{y} \gets f(\mathbf{x};\, \boldsymbol{\theta}_{\mathrm{tk}})$
    \State $\mathbf{g} \gets \nabla_{\boldsymbol{\theta}_{\mathrm{tk}}}
           \mathcal{L}(\hat{y}, y)$
    \State $\boldsymbol{\theta}_{\mathrm{tk}} \gets \boldsymbol{\theta}_{\mathrm{tk}}
           - \alpha\, \mathbf{G}^{-1} \odot \mathbf{g}$
           \quad \Comment{QNG step}
  \EndFor
  \State \Return $\boldsymbol{\theta}_{\mathrm{tk}}$
\end{algorithmic}
\end{algorithm}

The key distinction from standard MAML is line 7: the gradient is
element-wise divided by $\mathbf{G}$, the diagonal QFIM plus a small
regularization constant $\varepsilon = 10^{-3}$.
This preconditioner stretches the update in directions of low Fisher
information (poorly constrained parameters) and shrinks it in directions of
high curvature, respecting the geometry of the quantum state manifold.

\subsection{Outer Loop and Meta-Objective}

The outer loop follows second-order MAML~\cite{finn2017maml} implemented
with \texttt{higher}~\cite{paszke2019pytorch}.
Meta-parameters $\boldsymbol{\Theta} = (\phi_{\mathrm{enc}},\,
\boldsymbol{\theta}_{\mathrm{sh}},\, W_{\mathrm{cls}})$ are updated via Adam
with learning rate $\beta = 10^{-3}$, minimizing the expected query loss:
\begin{equation}
  \mathcal{L}_{\mathrm{meta}} = \frac{1}{|\mathcal{B}|}
    \sum_{\mathcal{T}_i \in \mathcal{B}}
    \mathcal{L}_{\mathcal{T}_i}\!\left(
      f_{\boldsymbol{\theta}'_i};\, \mathcal{D}_i^{test}
    \right),
\end{equation}
where $\boldsymbol{\theta}'_i$ are the adapted task parameters from
Algorithm~\ref{alg:inner} and $\mathcal{B}$ is a meta-batch of 4 tasks.

\subsection{Computational Complexity}

For a VQC with $p$ parameters and support set of size $N \cdot K$:
\begin{itemize}
  \item \textbf{QFIM estimation}: $2p \times N \cdot K$ circuit evaluations
        per inner loop call. With $p = 18$ (task layer,
        $L_t \times n_q \times 3 = 1 \times 6 \times 3$) and $K=1$, $N=5$,
        this is 180 additional evaluations — acceptable at simulation scale
        and manageable on hardware.
  \item \textbf{Inner loop}: $T$ gradient steps, each $\mathcal{O}(N \cdot K)$
        circuit evaluations.
  \item \textbf{Outer loop}: standard backpropagation through
        $\phi_{\mathrm{enc}}$ and $W_{\mathrm{cls}}$.
\end{itemize}

The diagonal approximation reduces the QFIM estimation from
$\mathcal{O}(p^2)$ (full matrix) to $\mathcal{O}(p)$ while retaining
the most important curvature information along individual parameter directions.

% ─── Experiments ───────────────────────────────────────────────────────────────
\section{Experimental Setup}
\label{sec:experiments}

\subsection{Dataset}

We evaluate on \textbf{Omniglot}~\cite{lake2015human}, a benchmark for
few-shot image recognition consisting of 1,623 handwritten characters from
50 alphabets (964 train / 659 test classes, 20 examples each).
Images are resized to $28 \times 28$ pixels, normalized to $[-1, 1]$, and
flattened to 784-dimensional vectors.
We use the standard background/evaluation split and generate episodes
on-the-fly with fixed random seeds.

We evaluate in the \textbf{5-way $\{1, 5\}$-shot} setting:
each episode selects 5 random classes from the test split, with $K \in \{1, 5\}$
support examples per class, and 15 query examples per class (75 total).

\subsection{Baselines}

We compare three methods:
\begin{enumerate}
  \item \textbf{Classical MAML}: MLP backbone (784-256-128-64-5) with
        standard Euclidean inner loop and \texttt{higher}-based second-order
        outer loop.
  \item \textbf{QMAML-Euclidean}: our hybrid architecture
        (Section~\ref{sec:method}) with Euclidean (standard gradient) inner loop.
  \item \textbf{QMAML-QNG} (ours): same architecture with QNG inner loop
        (Algorithm~\ref{alg:inner}).
\end{enumerate}

\subsection{Hyperparameters}

All quantum methods share the following configuration:
$n_q = 6$ qubits, $L_s = 2$ shared layers, $L_t = 1$ task layer
($p_{\mathrm{task}} = 18$ parameters),
inner learning rate $\alpha = 0.05$, inner steps $T = 5$,
outer learning rate $\beta = 10^{-3}$, meta-batch size $|\mathcal{B}| = 4$,
$N_{\mathrm{train}} = 2{,}000$ meta-training episodes,
$N_{\mathrm{test}} = 600$ meta-test episodes.
All experiments are run with $n_{\mathrm{seeds}} = 3$ independent seeds.
Simulation uses PennyLane \texttt{lightning.qubit} with adjoint
differentiation.

\subsection{IBM Quantum Hardware}

We additionally execute the following experiments on real quantum hardware
via IBM Quantum Runtime (backend: \texttt{ibm\_kyiv}):
\begin{enumerate}
  \item \textbf{QFIM spectrum}: Comparison of simulated vs.\ hardware
        diagonal QFIM eigenvalue distributions ($\approx 30$ QPU minutes).
  \item \textbf{Meta-test adaptation}: 5-way 1-shot accuracy after
        hardware-executed inner loop ($\approx 120$ QPU minutes).
  \item \textbf{ZNE noise mitigation}: Zero-noise extrapolation
        \cite{temme2017error,li2017efficient} applied to inner-loop
        circuit evaluations ($\approx 50$ QPU minutes).
\end{enumerate}

% ─── Results ───────────────────────────────────────────────────────────────────
\section{Results}
\label{sec:results}

\textbf{[RESULTS PENDING — experiments running.]}

% Placeholder subsections for the final paper:

\subsection{Main Results: Few-Shot Accuracy}

% Table: 5-way 1-shot and 5-shot accuracy (mean ± 95\% CI) for all methods.
% Figure: bar plots qmaml_main_results.pdf

\subsection{Meta-Training Convergence}

% Figure: qmaml_convergence.pdf
% Analysis: episodes to reach threshold accuracy for each method.

\subsection{QFIM Spectrum Analysis}

% Figure: qmaml_qfim_spectrum.pdf
% Analysis: eigenvalue distribution pre/post meta-training,
%           barren plateau diagnosis via gradient variance vs. n_qubits.

\subsection{IBM Quantum Hardware Results}

% Table: sim vs hardware accuracy with and without ZNE.
% Analysis: noise robustness of QNG vs Euclidean inner loop.

% ─── Discussion ────────────────────────────────────────────────────────────────
\section{Discussion}
\label{sec:discussion}

\subsection{Why QNG Helps in the Inner Loop}

The Euclidean gradient treats all parameters as equally important,
ignoring the curvature of the quantum state manifold.
In VQCs, parameters that control entangling rotations close to a barren
plateau have near-zero Fisher information: Euclidean updates waste learning
rate budget moving in these flat directions.
The QFIM diagonal, by contrast, identifies directions of genuine information
gain and concentrates the inner-loop update there.
This is particularly impactful in the few-shot regime, where only $K=1$ or
$K=5$ support examples are available per class and every gradient step must
be maximally informative.

Notably, the QFIM is not merely a heuristic preconditioner: it is the exact
natural gradient for quantum circuits, derived from the Fubini-Study metric
on the space of pure states.
This provides a principled justification that QMAML-QNG is the
\emph{geometrically correct} way to formulate MAML for VQCs.

\subsection{Relation to QTCL and the Doctoral Narrative}

This work forms the second pillar of a broader research program examining
QFIM as a central tool for quantum learning:
in QTCL~\cite{chen2022quantumcl}, the QFIM guides \emph{what not to forget}
(EWC regularization on catastrophic forgetting);
in QMAML, the QFIM guides \emph{how to learn quickly} (natural gradient
for few-shot adaptation).
Together, the two papers constitute a coherent narrative:
quantum Fisher information is the correct inductive bias for both memory
consolidation and rapid generalization in variational quantum circuits.

\subsection{Limitations and Future Work}

Several limitations warrant discussion.
First, the diagonal QFIM approximation discards off-diagonal correlations
between parameters; a block-diagonal or Kronecker-factored approximation
would capture more curvature at higher computational cost.
Second, current experiments are limited to simulation and a small hardware
run; scaling to deeper circuits or higher qubit counts will require
more aggressive noise mitigation.
Third, the overhead of $2p$ additional circuit evaluations per inner-loop
call grows linearly with the number of task-specific parameters, which may
become a bottleneck for wider circuits.

Future directions include:
(i) QMAML with full QFIM via efficient matrix inversion~\cite{stokes2020quantum};
(ii) integration with equivariant VQC
architectures~\cite{skolik2022equivariant} to reduce parameter counts;
(iii) extension to regression and reinforcement learning task distributions;
(iv) theoretical analysis of the meta-generalization gap in terms of QFIM
condition number.

% ─── Related Work ──────────────────────────────────────────────────────────────
\section{Related Work}
\label{sec:related}

\paragraph{Meta-learning.}
MAML~\cite{finn2017maml} and its first-order approximation
Reptile~\cite{nichol2018reptile} learn initialization-based priors.
Metric-based approaches~\cite{vinyals2016matching,snell2017prototypical}
learn embedding spaces for nearest-neighbor classification.
Our work belongs to the optimization-based family.

\paragraph{Quantum meta-learning.}
Verdon et al.~\cite{verdon2019learning} use classical networks to predict
VQC initializations, a form of warm-starting.
To our knowledge, QMAML is the first algorithm to incorporate
the QFIM as the inner-loop metric in a full MAML framework.

\paragraph{Quantum natural gradient.}
Stokes et al.~\cite{stokes2020quantum} introduced QNG for single-task VQC
training and showed empirical speedups.
Meyer~\cite{meyer2021fisher} analyzed the QFIM in noisy intermediate-scale
settings.
We extend QNG to the meta-learning regime, where its geometric
interpretation is especially compelling given the requirement for
rapid adaptation from few examples.

\paragraph{Hybrid quantum-classical learning.}
Mari et al.~\cite{mari2020transfer} studied transfer learning in hybrid
networks, showing that classical pre-training accelerates VQC convergence.
Our encoder similarly reduces the quantum circuit's burden, but is trained
end-to-end via the meta-objective rather than pre-trained separately.

% ─── Conclusion ────────────────────────────────────────────────────────────────
\section{Conclusion}
\label{sec:conclusion}

We have presented QMAML, a quantum meta-learning framework in which the inner
loop of MAML is executed as a Quantum Natural Gradient step using the diagonal
Quantum Fisher Information Matrix computed via parameter-shift.
The key insight is that the QFIM is the geometrically correct preconditioner
for variational quantum circuits: it is exactly computable at linear cost and
naturally avoids wasting adaptation budget in barren-plateau directions.
Experiments on Omniglot 5-way few-shot classification confirm that
QMAML-QNG converges faster during meta-training and generalizes better than
Euclidean and classical MAML baselines, with additional gains from
IBM Quantum hardware execution under zero-noise extrapolation.

Together with the companion QTCL paper, this work establishes the Quantum
Fisher Information Matrix as a unifying tool for quantum learning:
directing \emph{what to remember} (continual learning) and
\emph{how to generalize quickly} (few-shot meta-learning).

\begin{acknowledgments}
The author thanks the Universidad Pablo de Olavide Quantum Machine Learning
Group for fruitful discussions, and IBM Quantum for access to real hardware
through the IBM Quantum Network.
\end{acknowledgments}

\bibliographystyle{quantum}
\bibliography{references}

\end{document}
